Read in Gauge Theory Gravity: the metric $g$ below is the derived metric of the gravitational field $(\hb, \Omega)$ on flat spacetime, the manifold is the invertible region, and the tractor connection is the flat-space $\Cl(2,4)$ gauge field of Theorem \claimref{synthesis.gtg-tractor-gauge-field-reproduces-tractor-connection}; the index computation below is its reading in the splitting of $g$ (Fact \claimref{synthesis.imported-gtg-dictionary}).

(i) The curvature acts on $I$ by $F_{ab}\,I = [\nabla_a, \nabla_b]\,I = 0$, so at each point $F_{ab}$ lies in the stabiliser algebra of $I$. Since $Q(I) \neq 0$ the fibre splits as $\R I \oplus I^\perp$, and a bivector annihilating $I$ is a bivector of $I^\perp$ (write it in a basis adapted to the splitting: every term with an $I$ leg moves $I$). Parallel transport fixes $I$, hence preserves $I^\perp$ and the induced form on it, whose signature is $(1,4)$ or $(2,3)$ according to the sign of $\Lambda$ by Fact \claimref{dark-energy.imported-frozen-point-geometry}; so the holonomy lies in $SO(1,4)$ or $SO(2,3)$ and $\nabla$ restricts to a metric connection on $I^\perp$.
(ii) Lemma \claimref{dark-energy.stelle-west-vector-frozen-point}(i) gives $\ip{V}{I} = 0$ and $Q(V) = -3/\Lambda$; the subgroup of $SO(2,4)$ fixing both $I$ and $V$ acts on $\{I,V\}^\perp$, of signature $(2,4) - (1,0) - (0,1) = (1,3)$ for $\Lambda > 0$, where $I$ is positive and $V$ negative, and $(2,4) - (0,1) - (1,0) = (1,3)$ for $\Lambda < 0$, where they swap, and is its orthogonal group. Lemma \claimref{dark-energy.stelle-west-vector-frozen-point}(ii) gives $\nabla V = e$, Lemma \claimref{dark-energy.stelle-west-vector-frozen-point}(iii) gives the decomposition $D + e\wedge V/Q(V)$, and Lemma \claimref{dark-energy.stelle-west-vector-frozen-point}(iv) gives the quadratic term of the curvature: the bracket of the translational part with itself is $\tfrac12\cdot 2\cdot(\Lambda/3)\,e\wedge e$ in the Lorentz block, that is $\ell^{-2}e\wedge e$ with $\ell^2 = 3/\Lambda$, up to the normalisation of the bracket, while the mixed terms $[D, e\wedge V]$ carry the torsion $De$ into the translational block. This is the datum of Fact \claimref{dark-energy.imported-macdowell-mansouri-theorem-stelle-west} term by term.
(iii) On $I^\perp$ the fibre volume form is the six-dimensional volume form of $\mathcal{T}$ contracted with the unit normal $I/|I|$, and the Stelle-West contraction $\epsilon_{abcd}$ of Fact \claimref{dark-energy.imported-macdowell-mansouri-theorem-stelle-west} is that form contracted further with $V/\ell$; so $\epsilon_{abcd}F^{ab}F^{cd} = \epsilon(I, V, F, F)/(|I|\,\ell)$. Normalise with $|Q(I)|^{1/2}$ and $|Q(V)|^{1/2}$; since $Q(I)\,Q(V) = (\Lambda/3)(-3/\Lambda) = -1$ for either sign of $\Lambda$, the product of the two normalising factors is $1$ on both branches, and the Stelle-West prefactor $\ell^2 = 3/\Lambda = 1/Q(I)$ is all that remains, up to the fixed constant of the functional. The second form follows from $V = X - (3/\Lambda)I$ and the antisymmetry of $\epsilon$ in its first two slots. Slot count: six slots, two filled by $I$ and $V$, four by the two bivectors; the integrand is a scalar, which is why two vectors are needed and one is not enough.
(1) is inspection: every ingredient of $\epsilon(I, V, F, F)$, and of its value $\epsilon(I, X, F, F)$ at the canonical solution, is a section of the ambient bundle or its associated bundles, $F$ is polynomial in the connection, and no metric inverse occurs; the locators for the continuation of the bundle and the connection across the wall are those of Fact \claimref{dark-energy.imported-walls-sphere-theorem}.
(2) The pole is the rational-function statement $1/Q(I) = 3/\Lambda$, checked by the script; polynomiality of the integrand in $\Lambda$ holds because $I = (1, 0, \Lambda/6)$ and the connection \eqref{eq:conn} are polynomial in $\Lambda$. The expansion is Fact \claimref{dark-energy.imported-macdowell-mansouri-theorem-stelle-west} with $\ell^2 = 3/\Lambda$. The interior variation of the Gauss-Bonnet term vanishes: $\delta\int\epsilon_{abcd}R^{ab}\wedge R^{cd} = 2\int\epsilon_{abcd}\,D\delta\omega^{ab}\wedge R^{cd} = 2\int d\bigl(\epsilon_{abcd}\,\delta\omega^{ab}\wedge R^{cd}\bigr) - 2\int \epsilon_{abcd}\,\delta\omega^{ab}\wedge DR^{cd}$, the first a boundary term and the second zero by the Bianchi identity $DR = 0$. The remaining two terms are polynomial in $\Lambda$, and at $\Lambda = 0$ the cosmological term vanishes, leaving the Einstein-Cartan functional, whose only ingredients are $e$, $R$ and $\epsilon_{abcd}$: no inverse of $e$ is taken, so it is defined at degenerate $e$.
